% =====================================================================
%  THE COMPANION TO THE IGUSA SQUARE ROOT
%  Distilled from ~/igusa-cusp-form. Standalone document.
%  Author: Raeez Lorgat.
% =====================================================================
\documentclass[11pt]{article}
\usepackage[localtheorems]{raeez-math-template}

\providecommand{\cA}{\mathcal{A}}
\providecommand{\cC}{\mathcal{C}}
\providecommand{\cD}{\mathcal{D}}
\providecommand{\cF}{\mathcal{F}}
\providecommand{\cH}{\mathcal{H}}
\providecommand{\cL}{\mathcal{L}}
\providecommand{\cM}{\mathcal{M}}
\providecommand{\cO}{\mathcal{O}}
\providecommand{\cP}{\mathcal{P}}
\providecommand{\cW}{\mathcal{W}}
\providecommand{\cZ}{\mathcal{Z}}
\providecommand{\bC}{\mathbb{C}}
\providecommand{\bR}{\mathbb{R}}
\providecommand{\bQ}{\mathbb{Q}}
\providecommand{\bZ}{\mathbb{Z}}
\providecommand{\bH}{\mathbb{H}}
\providecommand{\bF}{\mathbb{F}}
\providecommand{\fg}{\mathfrak{g}}
\providecommand{\fD}{\mathfrak{D}}
\providecommand{\fA}{\mathfrak{A}}
\providecommand{\fM}{\mathfrak{M}}
\providecommand{\barB}{\overline{B}}
\providecommand{\MC}{\mathrm{MC}}
\providecommand{\Ran}{\operatorname{Ran}}
\providecommand{\Mbar}{\overline{\mathcal{M}}}
\providecommand{\FM}{\overline{\mathrm{FM}}}
\providecommand{\Vir}{\mathrm{Vir}}
\providecommand{\ChirHoch}{\operatorname{ChirHoch}}
\providecommand{\Zder}{Z^{\mathrm{der}}_{\mathrm{ch}}}
\providecommand{\kch}{\kappa_{\mathrm{ch}}}
\providecommand{\kBKM}{\kappa_{\mathrm{BKM}}}
\providecommand{\kcat}{\kappa_{\mathrm{cat}}}
\providecommand{\kfib}{\kappa_{\mathrm{fiber}}}
\providecommand{\SCchtop}{\mathsf{SC}^{\mathrm{ch,top}}}
\providecommand{\PhiFA}{\Phi^{\mathrm{FA}}}
\providecommand{\SpCh}{\mathrm{Sp}^{\mathrm{ch}}}
\providecommand{\CoHA}{\mathrm{CoHA}}
\providecommand{\Dfive}{\Delta_5}
\providecommand{\Phiten}{\Phi_{10}}
\providecommand{\Tr}{\operatorname{Tr}}
\providecommand{\Rez}{\operatorname{Res}}
\providecommand{\Conf}{\operatorname{Conf}}
\providecommand{\Pf}{\operatorname{Pf}}
\providecommand{\smult}{\operatorname{smult}}
\providecommand{\dend}{\operatorname{den}}
\providecommand{\Sp}{\mathrm{Sp}}
\providecommand{\SO}{\mathrm{SO}}
\providecommand{\HypPlan}{U}

\newtheorem{theorem}{Theorem}[section]
\newtheorem{proposition}[theorem]{Proposition}
\newtheorem{definition}[theorem]{Definition}
\newtheorem{conjecture}[theorem]{Conjecture}
\newtheorem{corollary}[theorem]{Corollary}
\newtheorem*{principle}{Principle}

\title{\bfseries The Igusa Square Root\\[0.4em]\Large The Companion}
\author{Raeez Lorgat}
\date{July 2026}

\begin{document}
\maketitle

\begin{abstract}
\noindent\itshape
The Igusa cusp form $\chi_{10}$ of weight ten on $\Sp_4(\bZ)$ has a
square root: the weight-five cusp form $\Dfive$ with the order-two
Maass multiplier $\nu_{\Dfive}$, and the square root, not the square,
is the fundamental object. The BPS partition function
$Z^{K3\times E}_{\mathrm{BPS}}=\Dfive^{-2}=\chi_{10}^{-1}$ is blind to
the sign of $\Dfive$; the Pfaffian --- the square root of the
determinant line --- retains the order-two orientation character, and
the manuscript's central conditional identity is
$\epsilon_o=\nu_{\Dfive}$. Five readings of $\Dfive$ are assembled:
the automorphic section, the denominator of a generalized
Borcherds--Kac--Moody superalgebra on a rank-three Lorentzian lattice,
a protected Pfaffian of a $K3\times E$ factorization datum, a
conjectural mirror discriminant, and the singular theta lift of the
$K3$ half-elliptic-genus $\phi_{0,1}$. The triangle
automorphic--denominator--theta is classical; the Pfaffian reading is
a recognition criterion relative to named finite data; the mirror
reading is open. This companion states the five views with exact
normalisations, the two unconditional new results, the honest ledger
of retained data, and the place of $\Dfive$ in the five-garden
programme: the terminal scalar shadow of the $K3\times E$ object.
\end{abstract}

\section{The seed}\label{sec:seed}

A partition function that is a perfect square asks for its square
root. On the trivial-canonical threefold $X=K3\times E$ the reduced
curve-counting partition function is
$Z^{X}_{\mathrm{BPS}}=\chi_{10}^{-1}$, and
$\chi_{10}=\Dfive^2$ where $\Dfive$ is the unique weight-five cusp
form on $\Sp_4(\bZ)$ with character --- the Maass multiplier
$\nu_{\Dfive}$, of order two, trivial on the theta group and equal to
$\det$ on the relevant reflection subgroup. Squaring kills exactly one
bit: an order-two twist that no scalar function of the period sees.
The manuscript's organising claim is that this bit is geometric. It is
the orientation of a moduli problem --- the $\bZ/2$ of a Pfaffian line
whose square is a determinant line --- and the five readings of
$\Dfive$ are five levels at which one and the same $K3\times E$ datum
$\fA_{K3\times E}\in E_3\text{-}\mathrm{HolFA}(K3\times E)$ presents
it:
\[
\textcircled{1}\ \text{automorphic section}\qquad
\textcircled{2}\ \text{BKM denominator}\qquad
\textcircled{3}\ \text{protected Pfaffian}\qquad
\textcircled{4}\ \text{mirror discriminant}\qquad
\textcircled{5}\ \text{theta lift.}
\]
The edges $\textcircled{1}\leftrightarrow\textcircled{2}$,
$\textcircled{2}\leftrightarrow\textcircled{5}$,
$\textcircled{1}\leftrightarrow\textcircled{5}$ are theorems of Igusa,
Borcherds and Gritsenko--Nikulin. The edge
$\textcircled{2}\leftrightarrow\textcircled{3}$ --- the
Pfaffian--Dirac identity --- is this manuscript's contribution, and it
is a \emph{recognition criterion}: a theorem relative to a named
finite package of retained data which the manuscript's own ledgers
record as not yet supplied. Reading $\textcircled{4}$ is a
conjecture. The five-fold thesis is deliberately not closed; what the
architecture supplies is the exact list of finite data that would
close it.

\section{The objects}\label{sec:objects}

\begin{definition}[The Jacobi seed]
$\phi_{0,1}$ is the weak Jacobi form of weight $0$ and index $1$,
normalised by the $K3$ elliptic genus
$Z_{K3}(\tau,z)=2\,\phi_{0,1}(\tau,z)$, with closed form
$\phi_{0,1}=4\sum_{i=2}^{4}\theta_i(\tau,z)^2/\theta_i(\tau,0)^2$ and
expansion $\phi_{0,1}=r^{-1}+10+r+O(q)$. Its coefficients $f(n,l)$
depend only on $4n-l^2$:
\[
f(0,0)=10,\quad f(0,\pm1)=1,\quad f(1,0)=108,\quad f(1,\pm1)=-64,
\quad f(1,\pm2)=10,\quad f(2,0)=808,\ \ldots
\]
with support $4n-l^2\ge-1$.
\end{definition}

\begin{definition}[The lattices]
With $\HypPlan$ the hyperbolic plane, the ambient rank-three lattice
is $\Lambda^{2,1}=\HypPlan\oplus\langle2\rangle$ with basis
$(f_2,f_3,f_{-2})$, $(f_2,f_{-2})=-1$, $(f_3,f_3)=2$. The type-II
sublattice $\Lambda^{2,1}_{II}$ (even coefficients on $f_{\pm2}$) has
index $4$ and discriminant $-32$. The simple type-II roots
\[
\delta_1=2f_2-f_3,\qquad \delta_2=2f_{-2}-f_3,\qquad \delta_3=f_3
\]
have Cartan matrix
$(\delta_i,\delta_j)=4\mathbf I_3-2\mathbf J_3=
\begin{psmallmatrix}2&-2&-2\\-2&2&-2\\-2&-2&2\end{psmallmatrix}$ and
Weyl vector $\rho=\tfrac12(\delta_1+\delta_2+\delta_3)$ with
$(\rho,\delta_i)=-1$. The off-diagonal entries $-2$ are infinite
Coxeter exponents: the reflection group
$W^{(2)}(\Lambda^{2,1}_{II})=\langle s_{\delta_1},s_{\delta_2},
s_{\delta_3}\rangle$ is the free product $\bZ/2*\bZ/2*\bZ/2$
--- the universal Coxeter group on three generators, with
abelianisation $(\bZ/2)^3$; there is no braid relation. The theta-lift
lattice is $\Lambda^{3,2}=\HypPlan\oplus\HypPlan\oplus\langle2\rangle$,
realised as $q_0^{\perp}$ inside $\wedge^2\bZ^4$ with the Pfaffian
pairing; the exceptional isogeny $\mathrm{PSp}_4(\bR)\simeq\SO(3,2)_+$
carries Siegel modular forms to orthogonal ones.
\end{definition}

\begin{definition}[Mukai--Gram dictionary]
The formal charge lattice of $X=S\times E$ is
$\Gamma^{\mathrm{form}}_X=\widetilde H(S,\bZ)\otimes
H^{\mathrm{even}}(E,\bZ)$, charges $(Q,P)$. The quadratic Gram map
$\Pi_X(Q,P)=\bigl(\tfrac{(Q,Q)}2,(Q,P),\tfrac{(P,P)}2\bigr)\in\bZ^3$
is not additive; its defect is a symmetric polarization $B$ forcing
the normal-ordered central extension
$\widehat\Gamma_X=\Gamma^{\mathrm{form}}_X\oplus_B\bZ^3$ on which the
corrected map is additive \textup{(the naive additive pushforward is
proved not to exist)}. For a curve class $\beta_h=s_1+hF$ with
$\chi(\cO_Y)=n$ and $d$ points, the D6--D2--D0 charge computes to
$\Pi_X=(h-1,\,n,\,d-1)$, and the root grading is
$\alpha(n,l,m)=2nf_2-lf_3+2mf_{-2}$ with
$\alpha(n,l,m)^2=-2(4nm-l^2)$.
\end{definition}

\begin{definition}[The Dirac--Igusa pro-object]
$\fD^{DI}_X$ is an inverse system over Harder--Narasimhan heights of
fourteen-component data: the $E_3$-prefactorization datum, a hybrid
factorization carrier on
$\Ran^{\mathrm{hyb}}(E)=\Ran(E)_{b=0}\sqcup\Ran^{\mathrm{wr}}(E)_{b>0}$,
the charge lattice with Gram projection, vanishing-cycle coefficients
with cosection-twisted reduced d-critical orientation \textup{(the
inflated $K3$ two-form supplies the cosection; Joyce--Upmeier square
roots supply the orientation candidates)}, Hall product and coproduct
on $H^\bullet_c(\fM^{\mathrm{red},+}(X),\Phi_W\otimes o)$, a
first-order Dirac block with Pfaffian line, a chiral Koszul source
coalgebra $C_X$ with cobar comparison towards
$U^{\mathrm{ch}}_E(\fg_{\Dfive}\otimes\cO_E)$, and a recognition map.
Two honesty clauses are part of the definition. The ``Dirac operator''
is algebraic: the first-order block at degree $\gamma$ is the
$2\times2$ skew matrix
$J_\gamma=\begin{psmallmatrix}0&1-x^\gamma\\-(1-x^\gamma)&0
\end{psmallmatrix}$ with
$\Pf(J_\gamma)=(1-x^\gamma)^{\operatorname{sdim}P_\gamma}$,
$x^\gamma=q^nr^ls^m$ --- no analytic operator on a constructed state
space is asserted. And the geometric Pfaffian line satisfies
$\cL_{\Pf}^{\otimes2}\simeq\cL_{\det}$: the square root of the
determinant of the cosection-reduced self-Ext complex, which is
exactly where the orientation bit lives.
\end{definition}

Three names discipline the one symbol: $\cD_X$ is the automorphic
section (view $\textcircled{1}$);
$\dend(\fg_{\Dfive})=64^{-1}\Dfive(2Z)$ is the denominator (view
$\textcircled{2}$); $\Pf_{\mathrm{prot}}(\fD^{DI}_X)$ is the protected
Pfaffian (view $\textcircled{3}$, conditional). The bare handle
$\Dfive$ never substitutes for any of the three.

\section{The theorems}\label{sec:theorems}

\begin{theorem}[View $\textcircled{1}$; proved]\label{thm:view1}
$\cD_X=\Dfive$ is the Borcherds--Gritsenko--Nikulin section of the
weight-five automorphic line on $\Sp_4(\bZ)\backslash\bH_2$ with
multiplier $\nu_{\Dfive}$: weight $5=f(0,0)/2$, divisor the simple
Humbert surface $\mathcal H_1$ \textup{(the diagonal locus)} with
multiplicity one, leading theta coefficient
$[q^{1/2}r^{1/2}s^{1/2}]\Dfive=2^6=64$, and in the chamber
$0<|s|\ll|q|\ll|r|^{-1}\ll1$ the Borcherds product
\[
\Dfive(Z)\;=\;64\,q^{1/2}r^{1/2}s^{1/2}
\prod_{(n,l,m)\in\Gamma_{\mathrm{eff}}}
\bigl(1-q^nr^ls^m\bigr)^{f(nm,l)} .
\]
The multiplier has order two, restricts to $\det$ on
$W^{(2)}(\Lambda^{2,1}_{II})$, and is trivial on chamber-permuting
type-I reflections.
\end{theorem}

\begin{theorem}[View $\textcircled{2}$; proved, Gritsenko--Nikulin]
\label{thm:view2}
The generalized Borcherds--Kac--Moody Lie superalgebra
$\fg_{\Dfive}$ on $\Lambda^{2,1}_{II}$ --- three even real simple
roots $\delta_i$, imaginary simple data with parity multiplicities
read off $\phi_{0,1}$ --- has Weyl--Kac--Borcherds denominator
\[
\dend(\fg_{\Dfive})
=\sum_{w\in W^{(2)}}\det(w)\,
w\Bigl(e^{-2\pi i(\rho,z)}-\!\!\sum_{a}m(a)\,e^{-2\pi i(\rho+a,z)}\Bigr)
=e^{-2\pi i(\rho,z)}\!\!\prod_{\alpha\in\mathcal R_+}\!\!
(1-e^{-2\pi i(\alpha,z)})^{\smult(\alpha)}
=\tfrac1{64}\Dfive(2Z),
\]
with signed root supermultiplicities
$\smult(\alpha(n,l,m))=f(nm,l)$ on the active support. The doubling
$Z\mapsto2Z$ records the index-four sublattice; $64$ is the theta
normalisation, not a doubling factor.
\end{theorem}

\begin{theorem}[View $\textcircled{5}$; proved, Borcherds]
\label{thm:view5}
The theta decomposition of $\phi_{0,1}$ is a weight $-1/2$
vector-valued form $h$ for the dual Weil representation of
$\langle2\rangle$; its regularized theta lift on $\Lambda^{3,2}$,
pulled through $\mathrm{PSp}_4(\bR)\simeq\SO(3,2)_+$, is
$\Theta(\phi^{\mathrm{Weil}}_{0,1})=\Dfive$, with weight
$c_{\phi}(0,0)/2$ and divisor
$\sum c_\phi(-\lambda^2/2,\lambda)\,\lambda^\perp$.
\end{theorem}

\begin{theorem}[The BPS scalar; proved elsewhere]\label{thm:bps}
For primitive classes on $K3\times E$ in the
Oberdieck--Pandharipande chamber, the reduced Donaldson--Thomas,
stable-pairs and Oberdieck--Pixton partition functions agree:
\[
Z^X_{\mathrm{DT}}=Z^X_{\mathrm{PT}}=Z^X_{\mathrm{OP}}
=-4096\,\Dfive^{-2},
\qquad
Z^{K3\times E}_{\mathrm{BPS}}:=(\Phiten^{\mathrm{un}})^{-1}
=\Dfive^{-2},
\]
where $4096=64^2$ is the squared theta-leading constant and the sign
is the Oberdieck--Pandharipande scalar branch. The arithmetic of the
square: $\Delta_{10}=\Dfive^2$ spans the Maass space, is the
Saito--Kurokawa lift of $\Delta E_6\in S_{18}(\mathrm{SL}_2(\bZ))$,
with spinor $L$-function
$L_{\mathrm{spin}}(s,\Delta_{10})=\zeta(s-8)\,\zeta(s-9)\,
L(s,\Delta E_6)$ --- a Ramanujan-violating CAP form, eigenvalues
$\lambda(2)=240$, $\lambda(3)=21960$.
\end{theorem}

\begin{theorem}[The two unconditional new results]\label{thm:new}
\textup{(i)} At the triple root
$\delta_{123}=\delta_1+\delta_2+\delta_3$ \textup{(norm $-6$)} the
parity-resolved multiplicities of $\fg_{\Dfive}$ are
\[
\dim_{\bar0}=29,\qquad \dim_{\bar1}=93,\qquad 29-93=-64=f(1,1),
\]
the $29$ even directions arising as $2$ independent Jacobi words
$[e_k,[e_i,e_j]]$ plus $3\times9$ brackets against the nine isotropic
imaginary-simple fibres \textup{(whence $\smult(a_{ij})=10=1+9$)}, the
$93$ odd directions as imaginary simple generators.
\textup{(ii)} A divisor, a leading coefficient, a scalar Borcherds
product, and its square do not determine a section of
$\cL^{5}\otimes T_\chi$ for a nontrivial order-two flat twist
$T_\chi$: order-two monodromy is invisible after squaring. This is
the precise sense in which $\Dfive$ carries strictly more than
$\chi_{10}$.
\end{theorem}

\begin{theorem}[Pfaffian--Dirac; conditional recognition]
\label{thm:pfdirac}
Granted the retained data $(D0)$ \textup{(degeneration limit of the
cosection-reduced orientation theory)}, $(O1)$ \textup{(strong reduced
orientation on the quotient self-Ext frame bundle)}, $(O1)^+$
\textup{(Weyl-equivariant transport along the three type-II
reflections)}, $(O2_{\mathrm{atlas}})$ \textup{(the local Pfaffian
wall atlas)} and the finite geometric Pfaffian datum
$(P_{\mathrm{fin}})$ on a cofinal Harder--Narasimhan system,
\[
\iota_{\mathrm{aut}}\bigl(\Pf_{\mathrm{prot}}(\fD^{DI}_X)\bigr)
=\cD_X=\Dfive,
\qquad
\epsilon_o=\nu_{\Dfive}\ \text{on}\ W^{(2)}(\Lambda^{2,1}_{II}).
\]
The orientation mechanism is a rank-one computation: at a Pfaffian
wall the odd crossing complex has $\Pf=\upsilon\,u$ and the
reflection flips its sign; since $W^{(2)}$ is a free product of three
involutions, a $\{\pm1\}$-character is fixed by its generator values,
where both $\epsilon_o$ and $\nu_{\Dfive}$ take $-1$. Honest scope:
the hypotheses of $(P_{\mathrm{fin}})$ encode the Fourier data of
$\Dfive$, so the identity is $q$-expansion rigidity relative to the
datum, not a construction; the manuscript's certificate ledgers record
every component of the retained package as not supplied; and the
$2^{12}=4096=64^2$ wall count is stipulated inside the atlas datum,
so its match with the squared theta constant carries no independent
evidence.
\end{theorem}

\begin{theorem}[Primitive recognition; conditional]\label{thm:recog}
Granted the ten-part finite compact-source datum $(S1)$--$(S10)$
\textup{(chain-level descent endpoint, Cartan matching, simple
representatives, Borcherds--Kac relations in the Hall bracket,
radical equality, no-extra-relations, generation, completed PBW,
exact triangular completion, and full parity dimensions
$\dim P_{\gamma,\bar p}=\operatorname{mult}^{\mathrm{GN}}_{\bar p}
(\alpha(\gamma))$ in both parities}\textup),
the protected primitive Lie superalgebra of the $K3\times E$ datum
recognises the Gritsenko--Nikulin algebra,
$\cP_X\cong\fg_{\Dfive}$, and the chiral Koszul source satisfies
$\Omega^{\mathrm{ch}}_E C_X\simeq
U^{\mathrm{ch}}_E(\fg_{\Dfive}\otimes\cO_E)$.
\end{theorem}

\begin{theorem}[Level-$\mathsf Z$ trace criterion; conditional]
\label{thm:trace}
Granted the retained trace datum $(Z_{\mathrm{HB}})$ and, as
established inputs, chiral Hochschild concentration on the Koszul
locus \textup{(Volume I)} and the chiral-Hochschild stratification
\textup{(Volume II)}, the bulk supertrace of the $K3\times E$ chart
satisfies
$\Tr^{\mathrm{bulk}}\bigl((-1)^F\bigr)_{\Zder(C_X)}=\Dfive^{-2}$, with
the Oberdieck--Pandharipande branch $-4096\,\Dfive^{-2}$ as the
separate scalar representative. The stronger level-$\mathsf A$
statement --- a gravity-line operator algebra acting on the boundary
with Pentagon trace $\Phiten^{\mathrm{un}}=\Dfive^2$ --- is open
\textup{(Volume II, second construction problem)}.
\end{theorem}

\begin{conjecture}[View $\textcircled{4}$: mirror discriminant]
\label{conj:mirror}
On the rank-three Mukai-invariant period component
$\Omega_{\widetilde M}/\Gamma_{\widetilde M}\simeq\bH_2/\Sp_4(\bZ)$
for $M=\HypPlan\oplus\langle-2\rangle$ \textup{(Dolgachev mirror
$\check M=\HypPlan\oplus E_7(-1)\oplus E_8(-1)$)}, the discriminant
section of the mirror period family equals $\Dfive^{-2}$; its polar
divisor is twice the reduced mirror discriminant, supported on the
Humbert divisor where $\Dfive$ vanishes. The theorem-level part is
the divisor identity $[\operatorname{div}\Dfive]=[\mathcal H_1]$; the
trivialisation, polar-order, and mirror-correspondence clauses are
open.
\end{conjecture}

\begin{conjecture}[Twined denominators]\label{conj:twined}
For $g$ in the Mathieu group $M_{24}$, the $g$-twined BKM denominators
satisfy $\dend(\fg^{(g)}_{\Dfive})
=\Theta(\phi^{(g),\mathrm{Weil}}_{0,1})|_{\bH_2}$; the classes
$7A,7B,15A,15B$ lie outside the CHL frame
\textup{(selected by $\varphi(N)\mid2$)} and their constant terms
need not be even --- the named arithmetic obstruction.
\end{conjecture}

\section{The computations}\label{sec:computations}

The five-coordinate $\kappa$-tuple of $K3\times E$, five distinct
constructions with no additive relation:
\[
\bigl(\kcat,\ \kch^{\mathrm{Hodge}},\ \kch^{\mathrm{Heis}},\
\kBKM,\ \kfib\bigr)=(0,\,0,\,3,\,5,\,24),
\]
where $\kcat=\chi(\cO_{K3})\chi(\cO_E)=0$ is K\"unneth-multiplicative,
$\kch^{\mathrm{Hodge}}=\sum_q(-1)^qh^{0,q}=0$,
$\kch^{\mathrm{Heis}}=3$ is the rank of the type-II Heisenberg--Cartan
after Gram projection, $\kBKM(\Dfive)=c_1(0)/2=5$ is the Borcherds
weight, and $\kfib=24$ is simultaneously $\chi_{\mathrm{top}}(K3)$ and
the Mukai rank --- numerically equal, semantically distinct. The
additive identity
$\kBKM=\kch^{\mathrm{Hodge}}+\chi(\cO_{\mathrm{fiber}})$ is false
\textup{(it predicts $2$; the value is $5$)}: the five coordinates are
incommensurable readings of one datum.

The universal Borcherds-weight identity is
$\kBKM(\Phi_N)=c_N(0)/2$, always naming the input denominator. On the
CHL ladder $N\in\{1,2,3,4,6\}$ the corrected weights are
\[
\mathrm{wt}=\bigl(5,\ 3,\ 2,\ \tfrac32,\ 1\bigr),
\qquad c_N(0)=(10,\ 6,\ 4,\ 3,\ 2),
\]
per Jatkar--Sen and Govindarajan--Krishna --- the weight at $N=4$ is
half-integral, and the once-recorded ladder $(5,4,3,2,1)$ is
retracted \textup{(its error: reading the frame-shape half-sum, the
weight of $\eta(\tau)^8\eta(2\tau)^8$, as a Jacobi constant term)}.
The Fake Monster form $\Phi_{12}$ carries $c_\Lambda(0)=24$,
$\kBKM=12$: one rule, distinct rows.

The first finite window of the denominator algebra, parity-resolved
($\mathrm{even}\,|\,\mathrm{odd}$): the isotropic doubles $2a_{ij}$
carry $10\,|\,0$; the triple root $\delta_{123}$ carries $29\,|\,93$;
the next window rows carry signed dimensions
$108$, $-513$, and at $2\delta_{123}$ superdimension $4016$ with
correction $m=-540$. The verification $[qrs]$-coefficient of the
monic product equals $93=64+29$.

Normalisation trap, stated once: $\Phiten^{\mathrm{un}}=\Dfive^2$
\textup{(theta-normalised)} while the Oberdieck--Pixton-monic form is
$\chi_{10}^{\mathrm{OP}}=D_5^2=4096^{-1}\Dfive^2$ with
$D_5=64^{-1}\Dfive$; the two differ by $2^{12}$, and every identity
above names its branch.

\section{The garden}\label{sec:garden}

One logarithmic $1$-form generates the programme. On a smooth complex
curve $X$, the form $\eta=d\log(z_1-z_2)$ is the unique
coordinate-free $1$-form with a residue along the diagonal of
$\Conf_2(X)$; the Arnold relation
$\eta_{12}\wedge\eta_{23}+\eta_{23}\wedge\eta_{31}
+\eta_{31}\wedge\eta_{12}=0$ generates all relations among its higher
companions. From $\eta$ one builds, for an augmented chiral algebra
$\cA$ on $X$, the ordered bar complex
$\barB^{\mathrm{ord}}_X(\cA)=T^c(s^{-1}\bar\cA)\otimes
\Omega^\bullet_{\log}(\FM_\bullet(X))$; its differential squares to
zero exactly when $\cA$ satisfies the Borcherds identity. Assembled
over $\Mbar_{g,n}$ with the prime-form propagator, the total
differential is a Maurer--Cartan element $\Theta_\cA=D_\cA-d_0$;
genus~$\ge1$ flatness fails by the curvature
$d^2_{\mathrm{fib}}=\kappa(\cA)\,\omega_g$, and the scalar $\kappa$
--- $\kappa(\Vir_c)=c/2$,
$\kappa(V_k(\fg))=\dim\fg\,(k+h^\vee)/2h^\vee$ --- is the central
charge read as modular curvature. The Verdier--Koszul dual obeys
conductor laws $c(\cA)+c(\cA^!)=K$ with $K(\Vir)=26$,
$K(\cW_N)=4N^3-2N-2$, $K(V_k(\fg))=2\dim\fg$ along
$k\mapsto-k-2h^\vee$; the landscape stratifies into archetype classes
$\mathsf G/\mathsf L/\mathsf C/\mathsf M$ of shadow depth
$2/3/4/\infty$ plus the critical Feigin--Frenkel stratum.

Every object sits on one chain: primitive factorization category
$\to$ chart algebra $A_b$ $\to$ bar shadow $\barB(A_b)$ $\to$ derived
chiral centre $\Zder(A_b)$ \textup{(the bulk)} $\to$ acting line
operators $\to$ scalar modular shadow. Identifying stages without the
licensing datum --- shadow $=$ object --- is the master error the
programme is disciplined against; this volume sits deliberately at
the terminal stage and refuses every unlicensed promotion: $\Dfive$
is not a Hilbert space, $Z_{\mathrm{BPS}}$ is not a path integral,
the denominator is not the bulk.

The five gardens are five readings of the chain. \emph{Volume I}
(\texttt{chiral-bar-cobar}) builds the engine on algebraic curves:
bar--cobar adjunction and inversion, quadratic Koszul recognition,
centre complementarity, the deformation--Hodge--graph traces
$F_g=\kappa\lambda_g^{\mathrm{FP}}+\delta F_g^{\mathrm{cross}}$,
Hochschild concentration, and the arithmetic shadow. \emph{Volume II}
(\texttt{chiral-bar-cobar-vol2}) reads the same object as physics:
bar $=$ BV/BRST, quantum master equation $=$ Maurer--Cartan, anomaly
$=$ curvature ($c-26\leftrightarrow\kappa_{\mathrm{tot}}$), the
two-coloured Swiss-cheese operad as bulk--boundary home, and
Universal Holography $\cA\mapsto(\partial=\cA,\
\mathrm{bulk}=\Zder(\cA))$, at $\Vir_{c=3\ell/2G_N}$ the boundary
reading of 3d gravity. \emph{Volume III}
(\texttt{calabi-yau-quantum-groups}) constructs the geometric source:
the two-stage functor $\Phi_d=\SpCh_{\Sigma_{d-1},C}\circ\PhiFA_d$
from Calabi--Yau $d$-categories to chiral algebras on a curve, the
cohomological Hall algebra $Y^+(X)$ with Drinfeld double $G(X)$ as
the level-three quantum group --- at $K3\times E$ the conditional
Hall--Borcherds recognition of $\fg_{\Dfive}$ that this volume's
criteria make precise. \emph{Mixed holomorphic--topological strings}
computes the local germ of the $d=2$ theory at $N$ Dirac branes,
$J(f)=\Tr f(\phi_1,\phi_2)$, with the holomorphic de Rham obstruction
$H^1(K3\times E,\cO)\neq0$ showing exactly why the local-to-global
route needs more than the formal stalk. \emph{This volume} is the
terminal scalar: the pentadic $\Dfive$, the one order-two bit that
survives every projection, and the finite list of data that would
promote it back to an operator.

The single question underneath: characterise the modular convolution
algebra $\fg^{\mathrm{mod}}_\cA$ as a geometric object. The proposed
answer --- endomorphisms of the Lagrangian embedding
$\cL_\cA\hookrightarrow\cM_{\mathrm{vac}}(\cA)$ in a $(-2)$-shifted
symplectic stack of vacua --- has Volume I computing its Taylor jets,
Volume II its local composition law, Volume III its geometric source,
and this volume its sharpest example: the $K3\times E$ datum whose
scalar shadow is $\Dfive$.

\section{The frontier}\label{sec:frontier}

The manuscript's architectural invention is the recognition-criterion
format: every conditional theorem names the exact finite data that
would close it, and the ledgers record what is supplied. The open
list, in order of depth:

\begin{enumerate}
\item The finite compact-source datum $(S1)$--$(S10)$ for primitive
recognition, and the finite Koszul comparison datum --- the operator-level
construction of $\fg_{\Dfive}$ from $K3\times E$ moduli.
\item The orientation package: $(D0_{\mathrm{HN}})$ degeneration
datum, $(O1_{\mathrm{quot}})$ quotient orientation,
Weyl transport $\mathcal T_W$, and the type-II wall atlas
$(O2_{\mathrm{atlas}})$ --- what makes $\epsilon_o$ a character
globally.
\item The finite geometric Pfaffian datum $(P_{\mathrm{fin}})$,
replacing $q$-expansion rigidity by construction.
\item The trace datum $(Z_{\mathrm{HB}})$ and the Hall--Borcherds
residual --- the four-part chain datum whose derived-centre trace is
$\Dfive^{-2}$.
\item The gravity-line operator algebra with Pentagon trace
$\Dfive^2$ \textup{(Volume II's second construction problem)} --- the
level-$\mathsf A$ completion.
\item The mirror clauses \textup{(M1)}--\textup{(M4)} of view
$\textcircled{4}$.
\item Twined denominators for all twenty-six $M_{24}$ classes, with
the parity obstruction at the non-CHL classes.
\end{enumerate}

What stands unconditionally: the classical triangle
$\textcircled{1}\textcircled{2}\textcircled{5}$ with every
normalisation pinned; the Oberdieck--Pandharipande scalar
$-4096\,\Dfive^{-2}$; the Saito--Kurokawa arithmetic of the square;
the parity-resolved window $29\,|\,93$ with $29-93=f(1,1)$; and the
monodromy proposition that makes the square root a genuine invariant.
The pentad is a programme; the square root is already a theorem.

\end{document}
